%% Results, discussion and conclusion. Tables are generated by paper/make_tables.py.

\begin{table}[t]
\centering\scriptsize
\caption{Data. Out-of-sample (OOS) bars follow the calibration period (volatility study, $h=24$
hourly, $h=5$ daily).}
\label{tab:data}
\resizebox{\linewidth}{!}{\begin{tabular}{lllllllrr}
\toprule
Set & Series & Ticker & Source & Freq. & Start & End & Bars & OOS bars \\
\midrule
Development & BTC & BTCUSDT & Binance & 1h & 2018-01-01 & 2026-09-30 & 76,544 & 53,581\\
 & ETH & ETHUSDT & Binance & 1h & 2018-01-01 & 2026-09-30 & 76,544 & 53,581\\
 & PAXG & PAXGUSDT & Binance & 1h & 2020-08-28 & 2026-09-30 & 53,354 & 37,348\\
 & S\&P 500 & \textasciicircum{}GSPC & Yahoo & 1d & 1927-12-30 & 2026-09-29 & 24,802 & 17,362\\
 & Nasdaq-100 & QQQ & Yahoo & 1d & 1999-03-10 & 2026-09-28 & 6,931 & 4,852\\
 & Gold & GLD & Yahoo & 1d & 2004-11-18 & 2026-09-28 & 5,498 & 3,849\\
 & Treasuries & TLT & Yahoo & 1d & 2002-07-30 & 2026-09-28 & 6,080 & 4,256\\
 & EUR/USD & EURUSD=X & Yahoo & 1d & 2003-12-01 & 2026-09-28 & 5,920 & 4,144\\
 & NIFTY 50 & \textasciicircum{}NSEI & Yahoo & 1d & 2007-09-17 & 2026-09-29 & 4,670 & 3,115\\
\midrule
Confirmatory & SOL & SOLUSDT & Binance & 1h & 2020-08-11 & 2026-09-30 & 53,775 & 37,643\\
 & BNB & BNBUSDT & Binance & 1h & 2018-01-01 & 2026-09-30 & 76,551 & 53,586\\
 & XRP & XRPUSDT & Binance & 1h & 2018-05-04 & 2026-09-30 & 73,625 & 51,538\\
 & ADA & ADAUSDT & Binance & 1h & 2018-04-17 & 2026-09-30 & 74,037 & 51,826\\
 & LTC & LTCUSDT & Binance & 1h & 2018-01-01 & 2026-09-30 & 76,551 & 53,586\\
 & DAX & \textasciicircum{}GDAXI & Yahoo & 1d & 1987-12-30 & 2026-09-29 & 9,799 & 6,860\\
 & Nikkei 225 & \textasciicircum{}N225 & Yahoo & 1d & 1965-01-05 & 2026-09-28 & 15,174 & 10,622\\
 & FTSE 100 & \textasciicircum{}FTSE & Yahoo & 1d & 1984-01-03 & 2026-09-29 & 10,797 & 7,558\\
 & Hang Seng & \textasciicircum{}HSI & Yahoo & 1d & 1986-12-31 & 2026-09-28 & 9,807 & 6,865\\
 & Russell 2000 & IWM & Yahoo & 1d & 2000-05-26 & 2026-09-29 & 6,624 & 4,637\\
 & Emerging mkts & EEM & Yahoo & 1d & 2003-04-14 & 2026-09-29 & 5,903 & 4,133\\
 & Silver & SLV & Yahoo & 1d & 2006-04-28 & 2026-09-29 & 5,137 & 3,582\\
 & Oil & USO & Yahoo & 1d & 2006-04-10 & 2026-09-29 & 5,150 & 3,595\\
 & GBP/USD & GBPUSD=X & Yahoo & 1d & 2003-12-01 & 2026-09-28 & 5,935 & 4,155\\
 & USD/JPY & JPY=X & Yahoo & 1d & 1996-10-30 & 2026-09-28 & 7,756 & 5,430\\
\bottomrule
\end{tabular}
}
\end{table}

%% ==================================================================
\section{Results}
\label{sec:results}

Table~\ref{tab:hyp} and Figure~\ref{fig:hyp} summarise the pre-registered tests. Every
development-set finding that was registered as a hypothesis replicated on the fifteen
confirmatory assets, with no deviation from the registered plan and no series excluded. The
subsections below report each part of the study, development and confirmatory sets side by side.

\begin{table}[t]
\centering\scriptsize
\caption{Pre-registered hypotheses. Development statistics are exploratory; confirmatory tests use
the registered code on fifteen new assets. Holm $p$ adjusts over H1--H8.}
\label{tab:hyp}
\resizebox{\linewidth}{!}{\begin{tabular}{l>{\raggedright\arraybackslash}p{5.2cm}>{\raggedright\arraybackslash}p{2.6cm}r>{\raggedright\arraybackslash}p{2.6cm}rr}
\toprule
& & \multicolumn{2}{c}{Development (exploratory)} & \multicolumn{3}{c}{Confirmatory (pre-registered)} \\
\cmidrule(lr){3-4}\cmidrule(lr){5-7}
ID & Hypothesis & Statistic & $p$ & Statistic & $p$ & Holm $p$ \\
\midrule
H1 & Analogue direction forecasts have zero mean IC across assets (null expected) & mean IC -0.0043; max DSR 0.58 & 0.562 & mean IC +0.0031; max DSR 0.46 & 0.603 & 1.000\\
H2 & Analogue+HAR QLIKE differs from HAR & median ratio 1.012 & 0.303 & median ratio 1.001 & 0.672 & 1.000\\
H3 & Analogue alone QLIKE differs from HAR & median ratio 1.103 & 0.001 & median ratio 1.087 & 0.000 & 0.000\\
H4 & Gaussian-kernel analogue beats kNN analogue (one-sided) & median ratio 0.967 & 0.000 & median ratio 0.953 & 0.000 & 0.000\\
H5 & Gradient boosting QLIKE differs from HAR & median ratio 1.522 & 0.000 & median ratio 1.423 & 0.000 & 0.000\\
H6 & Hawkes tail forecasts beat static Poisson (one-sided log-loss) & median ratio 0.930 & 0.000 & median ratio 0.847 & 0.000 & 0.000\\
H7 & Analogue-FHS blend FZ0 differs from FHS at alpha = 2.5\% & median diff -0.0003 & 0.910 & median diff -0.0005 & 0.934 & 1.000\\
H8 & Volatility targeting (Analogue+HAR) reduces max drawdown vs buy \& hold (one-sided sign test) & 9/9 improve & 0.002 & 15/15 improve & 0.000 & 0.000\\
\bottomrule
\end{tabular}
}
\end{table}

\begin{figure}[t]
\centering
\includegraphics[width=\linewidth]{fig_hypotheses.pdf}
\caption{Per-series log QLIKE ratios for H2--H4 (points) and their medians (bars). Negative values
favour the first-named model.}
\label{fig:hyp}
\end{figure}

\subsection{Direction: no predictability}
Table~\ref{tab:direction} reports the information coefficient (IC; Spearman correlation between
the forecast and the realised R-multiple of the trade), the Brier skill of the direction
probability, and out-of-sample Sharpe ratios. Across 24 assets the IC never exceeds 0.06 in
absolute value, the mean IC is \dirMeanICDev\ in the development set and \dirMeanICCon\ in the
confirmatory set, and H1 is confirmed ($p=0.60$). Richer embeddings ($m=4$) do not help. The
bar-for-bar replication of the original strategy (v1) and the walk-forward-selected v2 both trail
buy-and-hold on almost every asset with positive drift, and the best deflated Sharpe ratio over
the 18-configuration grid is \dirMaxDSRDev\ (development) and \dirMaxDSRCon\ (confirmatory), far below
the 0.95 needed to claim skill. Positive Brier skill appears only for Nasdaq-100 and gold (development) and oil (confirmatory);
all three had large changes in drift between the calibration and out-of-sample periods, which a
fixed climatological benchmark cannot track, and none has a significant IC.

The original system's claim that aggregating to higher frequencies acts as a ``low-pass filter''
receives only weak support: for Bitcoin the IC rises from 0.004 (15-minute bars) to 0.006
(hourly) and 0.02--0.035 (4-hour), but the 4-hour sample contains roughly a thousand independent
20-bar outcomes, and an IC of 0.03 is within sampling error.

\begin{table}[t]
\centering\scriptsize
\caption{Direction. IC: out-of-sample Spearman correlation between the analogue forecast and the
realised trade R-multiple. Sharpe ratios at 5\,bp per side. DSR: deflated Sharpe ratio of the
walk-forward-selected v2 over its 18-configuration grid.}
\label{tab:direction}
\resizebox{\linewidth}{!}{\begin{tabular}{lrrrrrrr}
\toprule
& \multicolumn{2}{c}{IC (long)} & Brier & \multicolumn{3}{c}{Sharpe ratio} & \\
\cmidrule(lr){2-3}\cmidrule(lr){5-7}
Asset & $m{=}1$ & $m{=}4$ & skill & B\&H & v1 & v2 WF & DSR \\
\midrule
\multicolumn{8}{l}{\textit{Development set}}\\
BTC & 0.006 & 0.028 & -0.028 & 0.84 & -0.44 & -0.12 & 0.01\\
ETH & 0.006 & 0.014 & -0.029 & 0.78 & -1.04 & -0.76 & 0.00\\
EUR/USD & -0.014 & -0.027 & -0.025 & -0.15 & -0.35 & -0.43 & 0.00\\
Gold & -0.040 & 0.012 & 0.077 & 0.71 & 0.34 & 0.40 & 0.42\\
NIFTY 50 & -0.058 & -0.030 & -0.051 & 0.62 & -0.00 & -0.64 & 0.00\\
Nasdaq-100 & -0.032 & -0.027 & 0.025 & 0.96 & 0.25 & 0.44 & 0.58\\
PAXG & -0.007 & 0.005 & -0.030 & 0.95 & -0.85 & -1.14 & 0.00\\
S\&P 500 & -0.006 & 0.005 & -0.057 & 0.55 & 0.10 & 0.08 & 0.32\\
Treasuries & 0.019 & -0.018 & -0.029 & -0.16 & 0.03 & -0.29 & 0.00\\
\midrule
\multicolumn{8}{l}{\textit{Confirmatory set}}\\
ADA & 0.011 & 0.018 & -0.031 & 0.67 & -0.51 & -0.43 & 0.00\\
BNB & -0.001 & 0.014 & -0.027 & 1.13 & -0.65 & -0.06 & 0.01\\
DAX & 0.002 & -0.007 & -0.047 & 0.37 & 0.17 & -0.12 & 0.01\\
Emerging mkts & 0.006 & -0.024 & -0.053 & 0.30 & -0.51 & -0.12 & 0.04\\
FTSE 100 & -0.001 & 0.021 & -0.068 & 0.27 & 0.14 & -0.05 & 0.01\\
GBP/USD & 0.013 & -0.023 & -0.026 & -0.12 & 0.21 & -0.08 & 0.06\\
Hang Seng & 0.005 & 0.001 & -0.078 & 0.25 & 0.12 & 0.02 & 0.06\\
LTC & 0.001 & 0.009 & -0.032 & 0.46 & -1.28 & -0.44 & 0.00\\
Nikkei 225 & 0.009 & 0.035 & -0.047 & 0.32 & 0.11 & -0.02 & 0.03\\
Oil & -0.034 & -0.052 & 0.113 & 0.34 & 0.03 & -0.11 & 0.02\\
Russell 2000 & -0.007 & 0.016 & -0.052 & 0.52 & -0.42 & 0.00 & 0.06\\
SOL & -0.010 & -0.001 & -0.037 & 0.82 & -1.33 & -0.92 & 0.00\\
Silver & 0.006 & -0.002 & -0.048 & 0.49 & -0.20 & -0.74 & 0.00\\
USD/JPY & -0.007 & 0.017 & -0.040 & 0.18 & -0.02 & -0.20 & 0.01\\
XRP & 0.004 & 0.027 & -0.028 & 0.81 & -0.95 & 0.66 & 0.46\\
\bottomrule
\end{tabular}
}
\end{table}

\subsection{Volatility: competitive in combination}
Table~\ref{tab:volsum} summarises 40 series--horizon pairs (full results in
Appendix~\ref{app:vol}). The analogue forecaster alone is 11\% worse than HAR on average and
significantly worse (H3). Combined with HAR it is statistically indistinguishable from HAR (H2;
median ratio 1.001 in the confirmatory set) and remains in the 90\% model confidence set in 23 of
25 confirmatory series, compared with 24 for HAR and 20 for GARCH(1,1); these two counts are
unchanged when bootstrap block lengths are chosen automatically (Appendix~\ref{app:rob}). This
replicates, prospectively and across asset classes, the finding of \citet{andrada2016} for the
S\&P 100: combining nearest-neighbour and model-based forecasts does not reliably improve
statistical accuracy. Gradient-boosted trees
given a superset of the same information are the worst model, with median QLIKE 42--52\% above
HAR, and are in the confidence set for 2 of 40 series (H5): with a few thousand effectively independent observations
the flexible model overfits the noisy volatility proxy, while the analogue estimator's
averaging over whole outcome windows limits its variance.

\begin{table}[t]
\centering\scriptsize
\caption{Volatility forecasting summary (QLIKE relative to HAR; lower is better).}
\label{tab:volsum}
\resizebox{\linewidth}{!}{\begin{tabular}{lrrrrrr}
\toprule
& EWMA & GARCH & HAR & GBM & Analogue & Analogue+HAR \\
\midrule
\multicolumn{7}{l}{\textit{Development set (15 series)}}\\
Mean QLIKE / HAR & 1.063 & 1.006 & 1.000 & 1.541 & 1.119 & 1.007\\
Median QLIKE / HAR & 1.049 & 1.003 & 1.000 & 1.522 & 1.103 & 1.012\\
In 90\% MCS & 11 & 14 & 14 & 0 & 8 & 14\\
Lowest QLIKE & 0 & 7 & 4 & 0 & 0 & 4\\
\midrule
\multicolumn{7}{l}{\textit{Confirmatory set (25 series)}}\\
Mean QLIKE / HAR & 1.077 & 1.007 & 1.000 & 1.566 & 1.115 & 1.002\\
Median QLIKE / HAR & 1.051 & 0.986 & 1.000 & 1.423 & 1.087 & 1.001\\
In 90\% MCS & 18 & 20 & 24 & 2 & 12 & 23\\
Lowest QLIKE & 1 & 14 & 2 & 0 & 1 & 7\\
\bottomrule
\end{tabular}
}
\end{table}

Figure~\ref{fig:cum} shows that the combination's advantage over HAR is not driven by a few
episodes: on Bitcoin the cumulative QLIKE gain rises steadily from 2024, while on the S\&P 500 the
combination is ahead of HAR for most of seven decades but loses ground in the 1987 and 2020 crashes. GARCH loses heavily to HAR on Bitcoin but leads on the S\&P
500, in line with the view that no single volatility model dominates \citep{hansen2005}.

\begin{figure}[t]
\centering
\includegraphics[width=\linewidth]{fig_cumulative_loss.pdf}
\caption{Cumulative QLIKE gain over HAR (development set). Rising lines beat HAR.}
\label{fig:cum}
\end{figure}

\paragraph{Ablations} Table~\ref{tab:abl} varies one design element at a time. The exclusion zone
and the full history each improve QLIKE by 2--3\%, confirming that independent outcomes and
many candidates matter. The Markov regime mixture has no measurable effect. The Gaussian kernel is
5\% better than $k$NN in both sets; this was found after the development protocol was fixed,
registered as H4, and confirmed ($p<10^{-6}$). Kernel weighting is therefore the recommended
default for future work.

\begin{table}[t]
\centering\scriptsize
\caption{Ablations: mean QLIKE relative to the full $k$NN analogue model ($<1$: the variant is
better).}
\label{tab:abl}
\resizebox{\linewidth}{!}{\begin{tabular}{lrrrrrrr}
\toprule
& \rotatebox{60}{m=1 (level only)} & \rotatebox{60}{m=2} & \rotatebox{60}{m=8} & \rotatebox{60}{no regime mixture} & \rotatebox{60}{no exclusion zone} & \rotatebox{60}{Gaussian kernel (bw 0.6)} & \rotatebox{60}{history capped at 5000}\\
\midrule
Development & 0.994 & 0.984 & 1.007 & 1.007 & 1.021 & 0.952 & 1.020\\
Confirmatory & 0.990 & 0.980 & 1.022 & 0.990 & 1.030 & 0.952 & 1.023\\
\bottomrule
\end{tabular}
}
\end{table}

\subsection{Tail events: clustering is volatility}
Table~\ref{tab:tail} evaluates forecasts of at least one raw jump in the next $h$ bars. The static
Poisson rule of the original system is \emph{worse than a constant forecast} in 11 of 15
development and 20 of 25 confirmatory series: a 1,000-bar average of past jump counts lags turning
points in risk. A Hawkes process fixes this (H6; median log-loss ratio 0.85 versus Poisson), but a
one-variable logistic regression on the EWMA variance is better still, with the lowest log-loss in
most series. Jump clustering is therefore largely volatility clustering, and a Hawkes intensity
adds little once the current variance is known. Figure~\ref{fig:rel} shows that the Poisson rule is
also miscalibrated: its highest-risk bin has a lower realised jump frequency than its middle bins.

\begin{table}[t]
\centering\scriptsize
\caption{Jump forecasts: probability of at least one raw jump ($|r|>3\bar\sigma$) in the next $h$ bars.}
\label{tab:tail}
\resizebox{\linewidth}{!}{\begin{tabular}{lrrrrr}
\toprule
& Constant & Poisson (v1) & Hawkes & Logistic-EWMA & Analogue \\
\midrule
\multicolumn{6}{l}{\textit{Development set}}\\
Mean log-loss / constant & 1.000 & 1.034 & 0.944 & 0.904 & 0.926\\
Mean AUC & 0.432 & 0.490 & 0.570 & 0.670 & 0.656\\
Best log-loss (series) & 0 & 0 & 0 & 12 & 3\\
Worse than constant & 0 & 11 & 1 & 0 & 2\\
\midrule
\multicolumn{6}{l}{\textit{Confirmatory set}}\\
Mean log-loss / constant & 1.000 & 1.061 & 0.901 & 0.893 & 0.911\\
Mean AUC & 0.437 & 0.502 & 0.613 & 0.710 & 0.679\\
Best log-loss (series) & 0 & 1 & 6 & 16 & 2\\
Worse than constant & 0 & 20 & 0 & 1 & 2\\
\bottomrule
\end{tabular}
}
\end{table}

\begin{figure}[t]
\centering
\includegraphics[width=0.5\linewidth]{fig_tail_reliability.pdf}
\caption{Reliability of jump-probability forecasts for Bitcoin (out-of-sample deciles).}
\label{fig:rel}
\end{figure}

\subsection{Value-at-Risk and Expected Shortfall: calibrated but not sharper}
Table~\ref{tab:var} compares one-day VaR/ES forecasts at 1\%, 2.5\% and 5\%. Historical simulation
and Gaussian GARCH fail coverage tests most often. The analogue-conditioned residual distribution
passes the Kupiec test as often as FHS (69\% of confirmatory cases for both) but ranks lower on the
FZ0 loss: 250 analogues per regime give a noisier estimate of a 1\% quantile than the full residual
history. Averaging the two (Blend) recovers FHS-level accuracy but does not improve on it (H7),
and FHS, Blend and GARCH-$t$ form the top group in both sets.

\begin{table}[t]
\centering\scriptsize
\caption{One-day VaR/ES. Pass rates are the share of series $\times$ $\alpha$ cases with $p>0.05$.}
\label{tab:var}
\resizebox{\linewidth}{!}{\begin{tabular}{lrrrrrr}
\toprule
& HS & GARCH-N & GARCH-t & FHS & Analogue & Blend \\
\midrule
\multicolumn{7}{l}{\textit{Development set}}\\
Mean $|$hit rate$/\alpha - 1|$ & 0.231 & 0.267 & 0.164 & 0.162 & 0.144 & 0.144\\
Kupiec pass rate & 0.63 & 0.59 & 0.78 & 0.85 & 0.89 & 0.89\\
Christoffersen CC pass rate & 0.41 & 0.56 & 0.67 & 0.74 & 0.81 & 0.74\\
Mean FZ0 rank (1 = best) & 5.81 & 3.85 & 1.89 & 2.30 & 4.74 & 2.41\\
In 90\% MCS (FZ0) & 0.52 & 0.85 & 1.00 & 1.00 & 0.63 & 1.00\\
\midrule
\multicolumn{7}{l}{\textit{Confirmatory set}}\\
Mean $|$hit rate$/\alpha - 1|$ & 0.243 & 0.338 & 0.228 & 0.176 & 0.184 & 0.162\\
Kupiec pass rate & 0.51 & 0.42 & 0.58 & 0.69 & 0.69 & 0.67\\
Christoffersen CC pass rate & 0.31 & 0.44 & 0.60 & 0.71 & 0.67 & 0.71\\
Mean FZ0 rank (1 = best) & 5.47 & 4.07 & 2.36 & 2.31 & 4.47 & 2.33\\
In 90\% MCS (FZ0) & 0.40 & 0.71 & 0.87 & 1.00 & 0.73 & 0.98\\
\bottomrule
\end{tabular}
}
\end{table}

\subsection{Economic value: lower drawdowns, not higher Sharpe ratios}
Scaling a long position by the Analogue+HAR volatility forecast reduced the maximum drawdown for
all 9 development and all 15 confirmatory assets (H8; median reduction 5.6 percentage points in the
confirmatory set) at the cost of slightly lower Sharpe ratios (median change $-0.015$; positive for
5 of 15 assets). This is the pattern documented by \citet{cederburg2020}: volatility management
reliably reduces risk but does not reliably improve risk-adjusted returns out of sample. With
dividend-adjusted prices for the four distribution-paying funds, Sharpe ratio levels rise but every
comparison keeps its sign (Appendix~\ref{app:rob}). The
Hawkes ``lock-out'' of the original system lowered the Sharpe ratio further for 9 of 15
confirmatory assets, and significantly so for Bitcoin and Ethereum in the development set, because the periods after
jumps include sharp rebounds (Table~\ref{tab:overlay}, Figure~\ref{fig:eq}).

\begin{table}[t]
\centering\scriptsize
\caption{Volatility-managed long exposure (no leverage, 5\,bp per unit turnover). $p$: paired
stationary-bootstrap test of the Sharpe difference with buy-and-hold.}
\label{tab:overlay}
\resizebox{\linewidth}{!}{\begin{tabular}{lrrrrrr}
\toprule
& \multicolumn{4}{c}{Sharpe ratio} & \multicolumn{2}{c}{Max drawdown (\%)} \\
\cmidrule(lr){2-5}\cmidrule(lr){6-7}
Asset & B\&H & Vol target & $p$ & + Hawkes gate & B\&H & Vol target \\
\midrule
\multicolumn{7}{l}{\textit{Development set}}\\
BTC & 0.84 & 0.78 & 0.30 & 0.34 & -77.2 & -75.9\\
ETH & 0.78 & 0.71 & 0.39 & 0.34 & -81.3 & -76.0\\
EUR/USD & -0.11 & -0.10 & 0.30 & -0.15 & -35.4 & -35.3\\
Gold & 0.44 & 0.47 & 0.32 & 0.47 & -45.6 & -42.7\\
NIFTY 50 & 0.72 & 0.72 & 0.83 & 0.72 & -38.4 & -31.2\\
Nasdaq-100 & 0.75 & 0.84 & 0.13 & 0.84 & -53.6 & -41.6\\
PAXG & 0.95 & 1.35 & 0.02 & 0.90 & -29.2 & -21.2\\
S\&P 500 & 0.55 & 0.58 & 0.19 & 0.60 & -56.8 & -49.3\\
Treasuries & -0.00 & -0.05 & 0.15 & -0.03 & -54.2 & -53.8\\
\midrule
\multicolumn{7}{l}{\textit{Confirmatory set}}\\
ADA & 0.67 & 0.42 & 0.00 & -0.23 & -95.5 & -95.0\\
BNB & 1.13 & 1.08 & 0.58 & 0.67 & -72.5 & -70.9\\
DAX & 0.37 & 0.38 & 0.86 & 0.41 & -72.7 & -62.7\\
Emerging mkts & 0.23 & 0.23 & 0.11 & 0.13 & -43.7 & -43.7\\
FTSE 100 & 0.27 & 0.22 & 0.25 & -0.01 & -52.6 & -46.6\\
GBP/USD & -0.09 & -0.10 & 0.31 & -0.05 & -37.5 & -37.4\\
Hang Seng & 0.25 & 0.23 & 0.63 & 0.10 & -65.2 & -58.3\\
LTC & 0.46 & 0.42 & 0.52 & 0.17 & -90.4 & -85.2\\
Nikkei 225 & 0.32 & 0.32 & 1.00 & 0.23 & -81.9 & -70.1\\
Oil & 0.09 & 0.08 & 0.85 & 0.08 & -94.6 & -85.2\\
Russell 2000 & 0.43 & 0.39 & 0.53 & 0.39 & -54.9 & -37.5\\
SOL & 0.82 & 0.83 & 0.93 & 0.19 & -82.7 & -77.2\\
Silver & 0.33 & 0.34 & 0.86 & 0.34 & -67.0 & -66.7\\
USD/JPY & 0.19 & 0.22 & 0.50 & 0.22 & -38.8 & -33.7\\
XRP & 0.81 & 0.60 & 0.10 & 0.22 & -84.8 & -75.1\\
\bottomrule
\end{tabular}
}
\end{table}

\begin{figure}[t]
\centering
\includegraphics[width=\linewidth]{fig_overlay_equity.pdf}
\caption{Growth of one unit, out-of-sample (development set).}
\label{fig:eq}
\end{figure}

\subsection{Interpretability and computation}
Every analogue forecast can be decomposed into its analogues. Figure~\ref{fig:explain} shows the
hour of the FTX collapse in November 2022: the ten nearest historical volatility shapes, what
followed them, and the forecast. Here the analogues predicted a partial reversion of volatility
that did not occur, which is the kind of error a user can see and reason about. The engine
forecasts about 30,000 bars per second with a capped history and 200,000 bars in about 90 seconds
with the full-history search on a 10-core laptop; the complete confirmatory study ran in under
half an hour.

\begin{figure}[t]
\centering
\includegraphics[width=\linewidth]{fig_analogue_explanation.pdf}
\caption{A forecast and its analogues: BTC, 8 November 2022, 18:00 UTC.}
\label{fig:explain}
\end{figure}

%% ==================================================================
\section{Discussion}

\paragraph{Where the predictability is} The analogue engine extracts no information about the
direction of prices but a substantial amount about their risk. This mirrors the broader
distinction between the unpredictability of returns \citep{fama1970} and the strong predictability
of volatility \citep{bollerslev1986,corsi2009}, and it holds across cryptocurrencies, equity
indices, commodities, bonds and currencies, at hourly and daily frequency.

\paragraph{Implications for the original system} Of the original system's components, the
directional analogue signal and the Poisson crash gate do not survive out-of-sample testing: the
former has no skill, and the latter is worse than a constant. The regime mixture, which replaced
an ad hoc adjustment, is harmless but adds nothing measurable. What survives is the core idea of
matching the present to the past, redirected from direction to risk, where it is competitive with
the standard models and yields well-calibrated tail forecasts with a transparent justification.

\paragraph{Why pre-registration mattered} Two development findings could have been overstated
without the confirmatory test. The analogue VaR was the best-calibrated model in development but
only tied with FHS on new assets; the kernel variant, found after the development design was
fixed, might have been dismissed as a lucky ablation but replicated with $p<10^{-6}$.
Pre-registration separated these cases at negligible cost.

\paragraph{Limitations} Realised variance is computed from bar returns rather than intraday
data, which adds noise to the target for daily series. The Analogue+HAR combination uses fixed
equal weights. The embedding uses only a single series; cross-asset analogues are a natural
extension. Transaction-cost assumptions are stylised, and the overlay ignores financing and cash
returns. The main tables use unadjusted closing prices, as registered; Appendix~\ref{app:rob} shows
that dividend adjustment does not change the conclusions.

%% ==================================================================
\section{Conclusion}
A pre-registered evaluation on 24 assets shows that analogue forecasting, a popular and
interpretable tool, has no directional value in financial markets but produces risk forecasts
that are competitive with HAR and GARCH, calibrated tail forecasts, and reliable drawdown
reductions. Its most useful design elements are an exclusion zone that keeps analogues
independent, the full history, and kernel weighting. We release the engine and the
pre-registration so that other analogue methods can be tested the same way.

\appendix
\section{Full volatility results}
\label{app:vol}
Bold entries are in the 90\% model confidence set.
\begin{center}\resizebox{0.9\linewidth}{!}{\begin{tabular}{lrrrrrr}
\toprule
Series (horizon) & EWMA & GARCH & HAR & GBM & Analogue & Analogue+HAR \\
\midrule
\multicolumn{7}{l}{\textit{Development set}}\\
BTC (h=24) & 1.101 & 1.107 & \textbf{1.000} & 1.270 & 1.041 & \textbf{0.992}\\
ETH (h=24) & 1.147 & \textbf{1.067} & \textbf{1.000} & 1.251 & \textbf{1.033} & \textbf{0.990}\\
EUR/USD (h=22) & \textbf{0.987} & \textbf{1.078} & \textbf{1.000} & 1.399 & \textbf{1.017} & \textbf{0.911}\\
EUR/USD (h=5) & \textbf{0.939} & \textbf{0.928} & 1.000 & 1.525 & 1.028 & 0.982\\
Gold (h=22) & \textbf{0.896} & \textbf{1.003} & \textbf{1.000} & 1.527 & \textbf{0.906} & \textbf{0.891}\\
Gold (h=5) & \textbf{1.015} & \textbf{0.969} & \textbf{1.000} & 1.625 & \textbf{1.103} & \textbf{1.012}\\
NIFTY 50 (h=22) & \textbf{1.044} & \textbf{0.894} & \textbf{1.000} & 2.016 & \textbf{1.272} & \textbf{1.071}\\
NIFTY 50 (h=5) & \textbf{1.073} & \textbf{1.004} & \textbf{1.000} & 1.574 & \textbf{1.106} & \textbf{1.023}\\
Nasdaq-100 (h=22) & \textbf{1.262} & \textbf{0.980} & \textbf{1.000} & 1.974 & 1.434 & \textbf{1.122}\\
Nasdaq-100 (h=5) & \textbf{1.022} & \textbf{0.942} & \textbf{1.000} & 2.139 & 1.084 & \textbf{1.018}\\
PAXG (h=24) & \textbf{1.109} & \textbf{1.087} & \textbf{1.000} & 1.308 & \textbf{1.172} & \textbf{1.047}\\
S\&P 500 (h=22) & 1.222 & \textbf{1.023} & \textbf{1.000} & 1.466 & 1.209 & \textbf{1.023}\\
S\&P 500 (h=5) & 1.075 & \textbf{0.984} & \textbf{1.000} & 1.152 & 1.056 & \textbf{0.992}\\
Treasuries (h=22) & \textbf{1.049} & \textbf{1.042} & \textbf{1.000} & 1.522 & \textbf{1.190} & \textbf{1.019}\\
Treasuries (h=5) & \textbf{1.010} & \textbf{0.975} & \textbf{1.000} & 1.368 & 1.133 & \textbf{1.012}\\
\addlinespace Mean ratio & 1.063 & 1.006 & 1.000 & 1.541 & 1.119 & 1.007\\
In 90\% MCS & 11/15 & 14/15 & 14/15 & 0/15 & 8/15 & 14/15\\
\midrule
\multicolumn{7}{l}{\textit{Confirmatory set}}\\
ADA (h=24) & 1.105 & 1.071 & \textbf{1.000} & 1.323 & \textbf{1.014} & \textbf{0.977}\\
BNB (h=24) & 1.151 & 1.162 & \textbf{1.000} & 1.242 & 1.070 & \textbf{1.001}\\
DAX (h=22) & 1.181 & \textbf{0.972} & \textbf{1.000} & 1.256 & 1.329 & 1.056\\
DAX (h=5) & 1.044 & \textbf{0.962} & 1.000 & 1.281 & 1.055 & 0.993\\
Emerging mkts (h=22) & \textbf{1.241} & \textbf{0.990} & \textbf{1.000} & \textbf{1.307} & \textbf{1.167} & \textbf{1.033}\\
Emerging mkts (h=5) & \textbf{1.051} & \textbf{0.984} & \textbf{1.000} & 1.358 & 1.098 & \textbf{1.012}\\
FTSE 100 (h=22) & \textbf{1.037} & \textbf{0.927} & \textbf{1.000} & 1.695 & 1.146 & \textbf{0.998}\\
FTSE 100 (h=5) & \textbf{0.993} & \textbf{0.936} & \textbf{1.000} & 1.410 & \textbf{1.017} & \textbf{0.980}\\
GBP/USD (h=22) & \textbf{0.983} & \textbf{0.991} & \textbf{1.000} & 1.624 & 1.108 & \textbf{1.009}\\
GBP/USD (h=5) & \textbf{1.000} & \textbf{0.982} & \textbf{1.000} & 1.645 & 1.087 & \textbf{1.011}\\
Hang Seng (h=22) & \textbf{1.017} & \textbf{0.950} & \textbf{1.000} & 1.756 & 1.231 & \textbf{1.034}\\
Hang Seng (h=5) & \textbf{1.011} & \textbf{0.986} & \textbf{1.000} & 1.423 & 1.091 & \textbf{1.008}\\
LTC (h=24) & 1.114 & 1.106 & \textbf{1.000} & 1.330 & 1.055 & \textbf{0.997}\\
Nikkei 225 (h=22) & \textbf{1.125} & \textbf{1.130} & \textbf{1.000} & 1.598 & \textbf{1.155} & \textbf{1.003}\\
Nikkei 225 (h=5) & 1.118 & 1.078 & \textbf{1.000} & 1.348 & 1.056 & \textbf{0.994}\\
Oil (h=22) & \textbf{1.362} & \textbf{0.986} & \textbf{1.000} & \textbf{2.214} & \textbf{1.486} & \textbf{1.164}\\
Oil (h=5) & \textbf{1.081} & \textbf{1.061} & \textbf{1.000} & 1.889 & \textbf{1.021} & \textbf{0.971}\\
Russell 2000 (h=22) & \textbf{1.333} & \textbf{0.924} & \textbf{1.000} & 1.808 & \textbf{1.634} & \textbf{1.172}\\
Russell 2000 (h=5) & \textbf{1.011} & \textbf{0.946} & \textbf{1.000} & 1.959 & 1.107 & \textbf{1.010}\\
SOL (h=24) & 1.105 & 1.110 & \textbf{1.000} & 1.313 & 1.053 & \textbf{0.996}\\
Silver (h=22) & \textbf{0.976} & \textbf{0.927} & \textbf{1.000} & 2.720 & \textbf{1.234} & \textbf{1.049}\\
Silver (h=5) & \textbf{0.979} & \textbf{0.970} & \textbf{1.000} & 1.600 & \textbf{1.002} & \textbf{0.960}\\
USD/JPY (h=22) & \textbf{0.658} & \textbf{1.058} & \textbf{1.000} & 1.432 & \textbf{0.650} & \textbf{0.691}\\
USD/JPY (h=5) & \textbf{1.012} & \textbf{0.977} & \textbf{1.000} & 1.375 & \textbf{1.021} & \textbf{0.978}\\
XRP (h=24) & \textbf{1.235} & \textbf{1.000} & \textbf{1.000} & 1.249 & \textbf{0.982} & \textbf{0.955}\\
\addlinespace Mean ratio & 1.077 & 1.007 & 1.000 & 1.566 & 1.115 & 1.002\\
In 90\% MCS & 18/25 & 20/25 & 24/25 & 2/25 & 12/25 & 23/25\\
\bottomrule
\end{tabular}
}\end{center}

\section{Robustness}
\label{app:rob}

\paragraph{Dividend-adjusted prices} The study uses Yahoo's unadjusted closing prices, as
registered. Price indices (S\&P 500, DAX and the like) have no adjusted series, and the gold,
silver and oil funds pay no distributions, so adjustment matters only for the Nasdaq-100,
long-Treasury, Russell 2000 and emerging-markets funds. Re-running the volatility and overlay
experiments on their dividend-adjusted prices (Table~\ref{tab:robadj}) moves QLIKE ratios by at
most 0.026. Sharpe ratio levels rise for both buy-and-hold and volatility targeting---most for
Treasuries, where coupons are the bulk of the return---but every comparison keeps its sign:
volatility targeting still lowers the maximum drawdown and still slightly lowers the Sharpe ratio
for three of the four funds.

\begin{table}[h]
\centering\scriptsize
\caption{Unadjusted versus dividend-adjusted prices (volatility: $h=5$ and $h=22$; overlay: $h=5$).}
\label{tab:robadj}
\resizebox{\linewidth}{!}{\begin{tabular}{llrrr}
\toprule
Series & Quantity & Unadjusted & Adjusted & Change \\
\midrule
Emerging mkts (h=22) & QLIKE / HAR: Analogue+HAR & 1.033 & 1.033 & -0.000\\
Emerging mkts (h=5) & QLIKE / HAR: Analogue+HAR & 1.012 & 1.007 & -0.005\\
Nasdaq-100 (h=22) & QLIKE / HAR: Analogue+HAR & 1.122 & 1.117 & -0.005\\
Nasdaq-100 (h=5) & QLIKE / HAR: Analogue+HAR & 1.018 & 1.012 & -0.006\\
Russell 2000 (h=22) & QLIKE / HAR: Analogue+HAR & 1.172 & 1.162 & -0.010\\
Russell 2000 (h=5) & QLIKE / HAR: Analogue+HAR & 1.010 & 1.006 & -0.004\\
Treasuries (h=22) & QLIKE / HAR: Analogue+HAR & 1.019 & 1.026 & +0.007\\
Treasuries (h=5) & QLIKE / HAR: Analogue+HAR & 1.012 & 1.010 & -0.002\\
Nasdaq-100 & Buy \& hold: Sharpe & 0.749 & 0.785 & +0.036\\
Nasdaq-100 & Buy \& hold: max DD & -0.536 & -0.534 & +0.001\\
Nasdaq-100 & Vol target: Analogue+HAR: Sharpe & 0.839 & 0.882 & +0.043\\
Nasdaq-100 & Vol target: Analogue+HAR: max DD & -0.416 & -0.409 & +0.007\\
Treasuries & Buy \& hold: Sharpe & -0.001 & 0.198 & +0.198\\
Treasuries & Buy \& hold: max DD & -0.542 & -0.484 & +0.058\\
Treasuries & Vol target: Analogue+HAR: Sharpe & -0.045 & 0.179 & +0.224\\
Treasuries & Vol target: Analogue+HAR: max DD & -0.538 & -0.446 & +0.092\\
Russell 2000 & Buy \& hold: Sharpe & 0.426 & 0.481 & +0.055\\
Russell 2000 & Buy \& hold: max DD & -0.549 & -0.543 & +0.006\\
Russell 2000 & Vol target: Analogue+HAR: Sharpe & 0.389 & 0.480 & +0.090\\
Russell 2000 & Vol target: Analogue+HAR: max DD & -0.375 & -0.352 & +0.023\\
Emerging mkts & Buy \& hold: Sharpe & 0.235 & 0.332 & +0.098\\
Emerging mkts & Buy \& hold: max DD & -0.437 & -0.398 & +0.039\\
Emerging mkts & Vol target: Analogue+HAR: Sharpe & 0.228 & 0.325 & +0.098\\
Emerging mkts & Vol target: Analogue+HAR: max DD & -0.437 & -0.393 & +0.045\\
\bottomrule
\end{tabular}
}
\end{table}

\paragraph{Bootstrap block lengths} Model confidence sets and Sharpe-difference tests use the
stationary bootstrap with fixed mean block lengths ($2h$; ten days for VaR/ES). Replacing them by
the automatic choice of \citet{politis2004}, with the correction of \citet{patton2009}, leaves the
headline counts unchanged---Analogue+HAR remains in the 90\% confidence set for 14 of 15
development and 23 of 25 confirmatory series, HAR for 14 and 24---and changes 19 of 672
confidence-set verdicts and 3 of 168 test verdicts (Table~\ref{tab:robblk}). The automatic lengths
are slightly more conservative, so the changes mostly admit additional models to a confidence set.

\begin{table}[h]
\centering\scriptsize
\caption{Fixed versus Politis--White bootstrap block lengths.}
\label{tab:robblk}
\resizebox{\linewidth}{!}{\begin{tabular}{llrlrrr}
\toprule
Set & Test & Cases & Verdict & Fixed & Politis--White & Changed \\
\midrule
Development & Volatility MCS (e2) & 90 & in MCS & 61 & 61 & 0\\
Development & VaR/ES MCS (e8) & 162 & in MCS & 135 & 135 & 2\\
Development & Overlay Sharpe test (e5) & 63 & significant & 4 & 4 & 0\\
Confirmatory & Volatility MCS (e2) & 150 & in MCS & 99 & 105 & 6\\
Confirmatory & VaR/ES MCS (e8) & 270 & in MCS & 211 & 220 & 11\\
Confirmatory & Overlay Sharpe test (e5) & 105 & significant & 14 & 13 & 3\\
\bottomrule
\end{tabular}
}
\end{table}

\paragraph{Kernel effective sample size} In kernel mode all anchors receive weight, and anchors
whose outcome windows overlap are not independent, so the \citet{kish1965} effective sample size
counted per anchor overstates the information in the analogue set. This affects the standard
error of the kernel forecast and its minimum-sample rule, not its mean, so H4, which compares
forecast accuracy, is unaffected. The released code offers an overlap-aware effective sample size
that counts each block of $h$ anchors as one observation; it is recommended for new work.

